% ---------------------------------------------------------------------
%  Hiện thực hóa use case (cuối mục Thiết kế chi tiết). Ma trận đầy đủ và
%  biểu đồ tuần tự từng use case nằm ở Phụ lục (backmatter/pl/e-hien-thuc-uc.tex).
% ---------------------------------------------------------------------
\subsection{Hiện thực hóa các use case trọng tâm}\label{ssec:hien-thuc-uc}

Các mục trước mô tả thiết kế theo chiều ngang: module, tầng và cơ chế dùng chung. Mục này đi theo chiều dọc, lần theo từng use case của Chương~\ref{chap:dac-ta} xuyên qua các tầng để kiểm chứng rằng mỗi bước, mỗi nhánh và mỗi quy tắc đều có phần tử thiết kế đảm nhận. Kỹ thuật sử dụng là hiện thực hóa use case (use case realization)~\cite{uml251}: use case được ánh xạ thành một cộng tác giữa các thành phần cụ thể và mô tả bằng biểu đồ tuần tự. Khác với biểu đồ hoạt động của Chương~\ref{chap:dac-ta} vốn coi hệ thống là hộp đen, các biểu đồ ở đây nhìn vào bên trong máy chủ, từ API qua dịch vụ ứng dụng và tầng miền tới PostgreSQL, Redis, River và dịch vụ bên ngoài.

Mọi use case theo cùng một khuôn. Use case được giao cho một dịch vụ ứng dụng (application service) của module sở hữu dữ liệu chính, gọi qua một API trong hợp đồng OpenAPI; riêng \uc{UC14} do tác vụ định kỳ của River khởi phát. Handler chỉ phân tích yêu cầu, xác thực và ánh xạ lỗi miền thành mã phản hồi. Dữ liệu vào được kiểm tra hợp lệ trước mọi truy vấn; mọi phép so sánh ngày nghiệp vụ dùng tham số \texttt{\$today} là ngày theo giờ Việt Nam do \texttt{platform/clock} cấp, không dùng \texttt{current\_date} của phiên CSDL vốn chạy theo UTC. Mỗi lệnh ghi chạy trong đúng một giao dịch do \texttt{TxManager} mở, với \texttt{statement\_timeout} và \texttt{lock\_timeout} 3 giây; khóa chết (40P01) hoặc quá thời gian chờ khóa được thử lại tối đa ba lần, sau đó trả 503 mã 5003. Với API có \texttt{Idempotency-Key}, câu lệnh đầu tiên của giao dịch là chèn khóa lũy đẳng ở trạng thái \texttt{in\_progress}; khóa đã hoàn tất thì phản hồi cũ được phát lại, khóa đang xử lý trả 409 mã 4002, cùng khóa khác nội dung trả 400 mã 2002, còn giao dịch lỗi bị hủy cùng khóa nên chỉ kết quả thành công được lưu (Hình~\ref{fig:seq-dong-bo-ngoai-tuyen}). Các khóa tiếp theo được lấy đúng thứ tự toàn cục ở mục~\ref{sssec:kt-fefo}: khóa tư vấn (advisory lock) theo không gian cho thao tác đổi thành viên, các dòng \texttt{space\_members} theo \texttt{user\_id}, rồi danh sách, món, mặt hàng, lô, và sau cùng bộ đếm \texttt{spaces.change\_seq}. Module khác tham gia qua cổng và dùng chung \texttt{pgx.Tx}. Mọi truy vấn đều lọc theo \texttt{space\_id}: người không phải thành viên của không gian trong đường dẫn, hoặc truy cập tài nguyên thuộc không gian khác, nhận 404 mã 4001 để không lộ sự tồn tại của tài nguyên (chống dò BOLA, \br{BR23}); thành viên thiếu vai trò cần thiết nhận 403 mã 3007 kèm \msg{MSG13}. Quy tắc nghiệp vụ được bảo vệ hai lớp: tầng miền kiểm tra trước để trả đúng mã quy tắc và thông điệp, còn ràng buộc CSDL (chỉ mục duy nhất từng phần, \texttt{CHECK}, khóa dòng) là hàng rào cuối khi có truy cập đồng thời. Tác dụng phụ ra ngoài giao dịch như đẩy thông báo, gửi thư, đánh chỉ mục hay xuất tệp được phát thành sự kiện miền, tức tác vụ River chèn trong cùng giao dịch (Hình~\ref{fig:seq-outbox-mua-hang}). Thực thể đồng bộ ghi \texttt{sync.space\_changes} ngay trước \texttt{COMMIT} và phát gợi ý thay đổi \texttt{space:\{sid\}} qua Redis sau commit (Hình~\ref{fig:seq-thoi-gian-thuc}); phép đọc trải nhiều bảng dùng giao dịch \texttt{REPEATABLE READ READ ONLY} để có ảnh chụp nhất quán. Xung đột phiên bản trả 409 mã 4002 kèm \msg{MSG35}, vượt giới hạn tần suất trả 429 mã 4290 kèm \msg{MSG36}. Khi một quy tắc buộc module thượng nguồn tác động lên dữ liệu của module hạ nguồn, module thượng nguồn định nghĩa giao diện móc (hook) và module hạ nguồn hiện thực nó lúc khởi động: \texttt{identity/api.ActivationHook} do household hiện thực để tạo không gian cá nhân, \texttt{household/api.MemberLeaveHook} do shopping hiện thực để trả mặt hàng đang giao về Cần mua (\br{BR07}), \texttt{recipe/api.RecipeDeleteHook} do mealplan hiện thực; nhờ vậy đồ thị phụ thuộc không có chu trình (mục~\ref{sec:kien-truc}).

Để giữ mạch trình bày của chương, phần chi tiết của việc hiện thực hóa được đặt ở Phụ lục~\ref{pl:hien-thuc-uc}. Phụ lục này gồm ma trận truy vết đủ 25 use case tới module, API, bảng dữ liệu, sự kiện và sơ đồ (Bảng~\ref{tab:c3-uc-thiet-ke}), cùng biểu đồ tuần tự phía máy chủ của các luồng chưa được minh họa ở mục~\ref{ssec:ky-thuat}: đăng ký và xác thực OTP (Hình~\ref{fig:seq-dang-ky-otp}), tham gia nhóm (Hình~\ref{fig:seq-tham-gia-nhom}), xóa thành viên, chuyển quyền và rời nhóm (Hình~\ref{fig:seq-roi-nhom-chuyen-quyen}), phân công người mua (Hình~\ref{fig:seq-phan-cong}), sinh danh sách từ kế hoạch (Hình~\ref{fig:seq-sinh-ds-ke-hoach}), thêm thực phẩm (Hình~\ref{fig:seq-them-thuc-pham}), gợi ý món (Hình~\ref{fig:seq-goi-y-mon}), lên lịch bữa ăn (Hình~\ref{fig:seq-lap-lich}), báo cáo và xuất tệp (Hình~\ref{fig:seq-bao-cao-xuat}), quản trị cấu hình, khóa tài khoản (Hình~\ref{fig:seq-quan-tri-cau-hinh}) và duyệt thực phẩm (Hình~\ref{fig:seq-duyet-thuc-pham}). Cùng với các hình ở mục~\ref{ssec:ky-thuat} và mục~\ref{sec:bao-mat}, mọi use case trọng tâm đều có biểu đồ hiện thực riêng.

Ma trận ở Phụ lục~\ref{pl:hien-thuc-uc} cho thấy ba dạng hiện thực: thao tác một thực thể trong một module (phần lớn gói tài khoản, công thức, quản trị); giao dịch trải nhiều module qua cổng (\uc{UC07}, \uc{UC11}, \uc{UC10} và \uc{UC19}), lý do chính để chọn một CSDL dùng chung; và phần xử lý bất đồng bộ qua River (\uc{UC05}, \uc{UC06}, \uc{UC08}, \uc{UC14}, \uc{UC22}, \uc{UC23}). Mười hai use case đặc tả đầy đủ ở Chương~\ref{chap:dac-ta} đều có biểu đồ tuần tự. Bốn use case \uc{UC09}, \uc{UC11}, \uc{UC14} và \uc{UC19} có biểu đồ riêng ở mục~\ref{ssec:ky-thuat}, còn \uc{UC13} dùng chung cơ chế FEFO của Hình~\ref{fig:seq-xac-nhan-nau}. Bảy use case còn lại là \uc{UC06}, \uc{UC08}, \uc{UC10}, \uc{UC11}, \uc{UC17}, \uc{UC18}, \uc{UC21}, cùng \uc{UC22}, được trình bày dưới đây. Ngoài nhóm đó, mục này hiện thực thêm những use case có luồng bảo mật hoặc đồng thời đáng chú ý: đăng ký, xác thực OTP và khôi phục mật khẩu (\uc{UC01} và \uc{UC03}); xóa thành viên, chuyển quyền và rời nhóm (\uc{UC07}); quản trị tài khoản, cấu hình và duyệt thực phẩm riêng (\uc{UC24}, \uc{UC23}, \uc{UC25}). Tìm kiếm và lọc (\uc{UC20}) có biểu đồ ở mục~\ref{sssec:kt-tim-kiem}.
